\section{Introduction}
Peer review is essential to scientific publishing as it evaluates the quality, reliability and reproducibility of research before publication. However, this process can be time-consuming and subject to variability and subjectivity between reviewers. Large language models (LLMs), such as Claude, ChatGPT and Gemini, could support this process by assisting with specific aspects of manuscript evaluation \cite{ebadi-2025}, but it remains unclear how reliable their assessments are compared with those of human reviewers \cite{suleiman-2024}.

Rather than asking whether LLMs could replace the entire peer review process, this study focuses on a narrower task: assessing the reporting quality of scientific papers with a standardized checklist, the  Checklist for Artificial Intelligence in Medical Imaging (CLAIM). Therefore, the main research question of this study is \textit{how closely do LLMs agree with human assessors when assessing the reporting quality of scientific papers with CLAIM?}. Additionally, the study examines whether LLMs tend to be stricter or more lenient than human assessors, whether agreement differs across models and which one is the most promising to support checklist-based evaluation in peer review.


\subsection{Literature review}
In medical imaging, transparency and reproducibility are fundamental for building trust in the clinical application of artificial intelligence (AI). For this reason, CLAIM was developed by experts as a standardized guideline for reporting AI studies in medical imaging. It is a structured checklist that covers the main sections of a scientific manuscript and addresses key aspects of the research process; including study design, data curation, model training and performance evaluation. Each item can be rated as Yes, No or Not Applicable; accommodating different study methodologies \cite{mongan-2020}. Following its revision in 2024, two additional items were incorporated, increasing the total number of checklist items from 42 to 44 \cite{tejani-2024}.

Because of these standardized, objective and verifiable criteria, recent editorial positions argue that checklist verification should become a mandatory audit during peer review \cite{zhong-2025}. However, applying these checklists consistently is challenging, as the process is time-consuming and varies between reviewers. 

In this context, LLMs have emerged as a potential tool for supporting peer review. Models such as ChatGPT, Claude and Gemini stand out in this context due to their capacity to reason over complex biomedical texts \cite{thirunavukarasu-2023}. Since human reviews can vary considerably due to subjectivity, LLMs might offer more objective assessments. However, their performance in unstructured manuscript evaluation has been inconsistent: studies have reported limited agreement between LLMs and human reviewers, including a potential leniency towards manuscripts with methodological flaws \cite{saad-2024},\cite{suleiman-2024}. In contrast, a blinded study found that, when the evaluation was structured AI-generated reviews were non-inferior to human ones \cite{zancanaro-2026}. This suggests that structured and standardized criteria may improve the consistency of LLM-assisted reviews. 

Recent studies suggest that this approach is feasible: models such as GPT-4 and Claude have been used to audit compliance with reporting guidelines such as CONSORT \cite{tsybulnik-2026}, and a study of neuroimaging papers using the COBIDAS guideline showed that AI-human agreement is possible but depends strongly on the LLM used \cite{pluimer-2026}.


\subsection{Rationale and justification}
Checklist-based assessment is central to systematic reviews, editorial screening and peer review, yet it is time-consuming and not required by all journals \cite{kocak-2024},\cite{kocak-2025}. This is particularly relevant in AI-based medical imaging, where incomplete reporting limits reproducibility and clinical translation. LLMs could make checklist audits faster and more scalable, but only if their judgments align with those of human assessors.

Evidence for CLAIM is still lacking. To our knowledge, no study has compared several LLMs with human assessors on CLAIM compliance scores under the same criteria, or examined whether LLMs are systematically stricter or more lenient than humans. By comparing three LLMs with human assessors, this study indicates which models, and in which CLAIM domains, come close enough to human scores to support reviewers and where human verification is still needed. In this way, it contributes to improving the rigor and efficiency of peer review, in line with the goals of meta-research: studying and improving how research itself is evaluated in the era of AI. 

  

